\documentclass[11pt]{article}

\usepackage[a4paper,margin=1in]{geometry}
\usepackage{amsmath,amssymb,amsthm,mathtools}
\usepackage{physics}
\usepackage{bm}
\usepackage{hyperref}
\usepackage{enumitem}

\title{From Two-Point Classical Action to Hamilton--Jacobi and Madelung Theory:\\
Hidden Boundary Data, Caustics, and the Breakdown of Gradient Velocity Fields}

\author{}
\date{}

\begin{document}

\maketitle

\begin{abstract}
The Hamilton--Jacobi equation is usually presented as a nonlinear first-order partial differential equation for a scalar field $S(q,t)$. However, the fundamental object of classical mechanics is not a one-point function but the two-point classical action
\[
S(q_0,t_0;q,t),
\]
defined as the extremal action between endpoints. The apparent reduction from a two-point object to a one-point field hides essential geometric and dynamical information. In particular, the Hamilton--Jacobi formulation is valid only while the evolving Lagrangian manifold remains a graph over configuration space. After caustic formation, the velocity field ceases to be globally representable as a gradient, and the single-valued phase function breaks down. We analyze the implications of this fact for Hamilton--Jacobi theory, Madelung hydrodynamics, and stochastic formulations of quantum mechanics.
\end{abstract}

\tableofcontents

\section{Introduction}

The Hamilton--Jacobi (HJ) equation occupies a central place in classical mechanics, geometrical optics, semiclassical quantum mechanics, and stochastic formulations of quantum theory. It is usually written as
\begin{equation}
\partial_t S + H(q,\nabla S,t)=0,
\label{HJ}
\end{equation}
where $S(q,t)$ is interpreted as a scalar field on configuration spacetime.

This formulation suggests that the theory is an ordinary initial value problem for a scalar field. However, this interpretation is misleading. The fundamental classical object is not the one-point function $S(q,t)$ but the two-point action
\begin{equation}
S(q_0,t_0;q,t),
\end{equation}
defined as the extremal value of the action functional between fixed endpoints.

The purpose of this article is to clarify:
\begin{enumerate}[label=(\roman*)]
\item how the HJ equation emerges from the two-point action;
\item why the apparent one-point formulation hides essential boundary information;
\item why the velocity field ceases to be globally representable as a gradient after caustic formation;
\item why Madelung hydrodynamics and related stochastic formulations inherit these difficulties.
\end{enumerate}

\section{The Classical Action as a Two-Point Function}

Consider a classical system with Lagrangian
\[
L(q,\dot q,t).
\]
The Hamilton principal function is defined by
\begin{equation}
S(q_0,t_0;q,t)
=
\operatorname*{ext}_{q(\tau)}
\int_{t_0}^{t} L(q,\dot q,\tau)\,d\tau,
\label{principal}
\end{equation}
where the extremization is performed over all paths satisfying
\[
q(t_0)=q_0,
\qquad
q(t)=q.
\]

Thus the classical action is naturally a two-point function on configuration spacetime.

The endpoint variation formulas yield
\begin{equation}
p_i(t)=\frac{\partial S}{\partial q^i},
\qquad
p_i(t_0)=-\frac{\partial S}{\partial q_0^i}.
\label{endpoint}
\end{equation}

Hence the two-point action generates the classical canonical transformation.

\section{Emergence of the Hamilton--Jacobi Equation}

Fix the initial endpoint $(q_0,t_0)$ and define
\begin{equation}
\mathcal S(q,t)
:=
S(q_0,t_0;q,t).
\end{equation}

Now vary only the final endpoint. A small variation gives
\begin{equation}
\delta S
=
p_i\,\delta q^i - H\,\delta t.
\end{equation}

Therefore,
\begin{equation}
dS = p_i\,dq^i - H\,dt.
\end{equation}

Identifying
\[
p_i=\frac{\partial S}{\partial q^i},
\]
one obtains
\begin{equation}
\frac{\partial S}{\partial t}
+
H\left(q,\frac{\partial S}{\partial q},t\right)
=
0.
\end{equation}

Thus the HJ equation appears only after one endpoint has been frozen.

\section{Hidden Boundary Data}

The reduction from the two-point action to the one-point field hides the initial endpoint inside the boundary conditions.

Indeed,
\[
S(q,t)=S(q_0,t_0;q,t)
\]
must satisfy singular initial data:
\begin{equation}
S(q,t_0)=
\begin{cases}
0, & q=q_0,\\
+\infty, & q\neq q_0.
\end{cases}
\end{equation}

Therefore the HJ equation alone does not define the dynamics. The choice of admissible solutions depends on hidden global information inherited from the original two-point action.

\section{Lagrangian Manifolds}

The geometric meaning of the HJ equation becomes clearer in phase space.

Given a function $S(q,t)$, define
\begin{equation}
p=\nabla S(q,t).
\end{equation}

This determines a submanifold
\begin{equation}
\Lambda_t
=
\{(q,p):p=\nabla S(q,t)\}
\subset T^*Q.
\end{equation}

Initially, one has a graph
\begin{equation}
\Lambda_0=\{(q,\nabla S_0(q))\}.
\end{equation}

Hamiltonian evolution transports this manifold:
\begin{equation}
\Lambda_t=\Phi_t(\Lambda_0),
\end{equation}
where $\Phi_t$ is the Hamiltonian flow.

The HJ formulation remains valid only while $\Lambda_t$ remains a graph over configuration space.

\section{Caustics and Breakdown of the Gradient Representation}

As the Hamiltonian flow evolves, the Lagrangian manifold generically folds in phase space. The projection
\[
\pi:\Lambda_t\to Q
\]
eventually ceases to be one-to-one.

At such points:
\begin{itemize}
\item several classical trajectories reach the same configuration point;
\item several momenta correspond to the same position;
\item the action becomes multivalued;
\item no single velocity field exists.
\end{itemize}

Instead of a single branch
\[
p=\nabla S(q,t),
\]
one obtains several branches:
\begin{equation}
p_1(q,t),\quad p_2(q,t),\quad \dots
\end{equation}

Each branch may locally be represented as a gradient,
\[
p_i=\nabla S_i,
\]
but globally no single-valued $S(q,t)$ exists.

\subsection{Physical Interpretation}

Suppose particles initially move to the right. A sufficiently strong potential may reverse some trajectories in finite time. Then at later times a single position may simultaneously contain:
\begin{itemize}
\item trajectories moving left;
\item trajectories moving right.
\end{itemize}

Therefore the velocity field is no longer single-valued.

This phenomenon is generic and corresponds mathematically to the intersection of characteristics in first-order PDE theory.

\section{Relation to Burgers Equation}

The HJ equation is closely related to Burgers-type equations.

For the free particle Hamiltonian
\[
H=\frac{p^2}{2m},
\]
one has
\begin{equation}
\partial_t S + \frac{(\nabla S)^2}{2m}=0.
\end{equation}

Defining
\[
v=\frac{\nabla S}{m},
\]
gives
\begin{equation}
\partial_t v + (v\cdot\nabla)v=0.
\end{equation}

This equation develops shocks through characteristic crossing. Thus the breakdown of the gradient representation is not exceptional but generic.

\section{Madelung Hydrodynamics}

The Madelung representation writes the wavefunction as
\begin{equation}
\psi=\sqrt{\rho}\,e^{iS/\hbar}.
\end{equation}

Substitution into the Schr\"odinger equation yields:
\begin{equation}
\partial_t \rho
+
\nabla\cdot
\left(
\rho\frac{\nabla S}{m}
\right)
=
0,
\label{continuity}
\end{equation}
and
\begin{equation}
\partial_t S
+
\frac{(\nabla S)^2}{2m}
+
V
-
\frac{\hbar^2}{2m}
\frac{\Delta\sqrt{\rho}}{\sqrt{\rho}}
=
0.
\label{madelungHJ}
\end{equation}

At first sight this resembles ordinary hydrodynamics. However, the velocity field
\[
v=\frac{\nabla S}{m}
\]
is highly constrained.

Away from singularities,
\begin{equation}
\nabla\times v=0.
\end{equation}

Thus Madelung fluids are not generic fluids but irrotational potential flows.

\section{Multivalued Phases and Interference}

After caustic formation, semiclassical quantum mechanics produces
\begin{equation}
\psi
\sim
\sum_j
A_j e^{iS_j/\hbar}.
\end{equation}

Thus quantum mechanics naturally accommodates multiple classical branches simultaneously.

The hydrodynamic velocity becomes
\begin{equation}
v
=
\frac{\hbar}{m}
\operatorname{Im}
\frac{\nabla\psi}{\psi},
\end{equation}
which generally is not equal to the gradient of a single classical action branch.

Near nodes where
\[
\psi=0,
\]
the velocity field becomes singular and vortices appear.

\section{Global Constraints and Quantization}

Single-valuedness of the wavefunction implies
\begin{equation}
\oint \nabla S\cdot dl
=
2\pi n\hbar.
\end{equation}

Thus the Madelung variables are globally constrained.

Without imposing these quantization conditions, the Madelung system admits more solutions than Schr\"odinger theory.

This issue is closely related to the Wallstrom objection in stochastic mechanics.

\section{Consequences for Nelson Mechanics}

Nelson stochastic mechanics introduces a current velocity
\[
v=\frac{\nabla S}{m}
\]
together with stochastic diffusion.

However, unless additional global constraints are imposed, the resulting theory generally possesses too many states compared with quantum mechanics.

The underlying issue is again the hidden geometric structure inherited from:
\begin{itemize}
\item the original two-point action;
\item the multivalued semiclassical structure;
\item global phase coherence.
\end{itemize}

\section{Conclusion}

The Hamilton--Jacobi equation does not truly replace a two-point dynamical object by a one-point field. The apparent scalar field $S(q,t)$ is only a local representation of an evolving Lagrangian manifold in phase space.

The gradient representation
\[
p=\nabla S
\]
fails generically after caustic formation, when the phase-space manifold folds and multiple classical trajectories reach the same configuration point.

Madelung hydrodynamics inherits these limitations. The velocity field is not globally well-defined, and the true quantum object remains the wavefunction itself rather than a classical fluid velocity field.

Thus the hidden two-point structure of classical action survives in:
\begin{itemize}
\item the geometry of Lagrangian manifolds,
\item multivalued semiclassical phases,
\item quantization conditions,
\item and the global topology of wave mechanics.
\end{itemize}

\begin{thebibliography}{99}

\bibitem{Arnold}
V. I. Arnold,
\textit{Mathematical Methods of Classical Mechanics},
Springer.

\bibitem{Maslov}
V. P. Maslov and M. V. Fedoriuk,
\textit{Semi-Classical Approximation in Quantum Mechanics},
Reidel.

\bibitem{Guillemin}
V. Guillemin and S. Sternberg,
\textit{Symplectic Techniques in Physics},
Cambridge University Press.

\bibitem{Madelung}
E. Madelung,
``Quantentheorie in hydrodynamischer Form,''
\textit{Zeitschrift f\"ur Physik} \textbf{40}, 322--326 (1927).

\bibitem{Nelson}
E. Nelson,
\textit{Quantum Fluctuations},
Princeton University Press.

\bibitem{Wallstrom}
T. C. Wallstrom,
``Inequivalence between the Schr\"odinger equation and the Madelung hydrodynamic equations,''
\textit{Phys. Rev. A} \textbf{49}, 1613 (1994).

\bibitem{Evans}
L. C. Evans,
\textit{Partial Differential Equations},
AMS.

\bibitem{Courant}
R. Courant and D. Hilbert,
\textit{Methods of Mathematical Physics},
Wiley.

\end{thebibliography}

\end{document}
